\documentclass[11pt]{article}

\usepackage[margin=1in]{geometry}
\usepackage[T2A]{fontenc}
\usepackage[utf8]{inputenc}
\usepackage[english,russian]{babel}
\usepackage{amsmath}
\usepackage{amssymb}
\usepackage{booktabs}
\usepackage{hyperref}
\usepackage{graphicx}
\usepackage{array}
\usepackage{authblk}

\title{Поиск направленных градиентов в популяциях галактик и наблюдаемых величинах квазаров высокого красного смещения по данным DESI DR1}

\author{I. Mikhailov}
\affil{Независимый исследователь, Москва, Россия}

\date{28 мая 2026}

\begin{document}

\maketitle

\begin{abstract}
Мы ищем крупномасштабные направленные градиенты в каталогах галактик и квазаров DESI Data Release 1 (DR1). Вопрос здесь эмпирический: остается ли зависимость плотности трассеров или их свойств от направления на небе при фиксированном красном смещении после учета геометрии обзора, эффектов отбора и известных наблюдательных систематик? Для анализа плотности мы используем clustering-каталоги DESI DR1 Large-Scale Structure и соответствующие random-каталоги, а для популяционных и эмиссионных характеристик — value-added каталог DESI DR1 FastSpecFit Iron. Набор проверок включает карту остатков LRG \texttt{DN4000\_MODEL} при $0.4 \le z < 0.6$, карты плотности QSO высокого красного смещения при $1.5 \le z < 3.5$, остатки эквивалентных ширин эмиссионных линий QSO и предварительно зафиксированное семейство отношений линий, построенное из \texttt{CIV\_1549\_EW}, \texttt{CIII\_1908\_EW} и \texttt{MGII\_2796\_EW}. В анализе используются HEALPix-карты неба, DESI random-каталоги там, где они нужны, регрессия по красному смещению и служебным переменным, внешние шаблоны систематик imaging-данных и пространственные блочные тесты нулевой гипотезы для амплитуд диполя. В контролируемом наблюдательном семействе минимальное блочное p-value равно $p=0.409$. В предварительно зафиксированном семействе отношений линий QSO минимальное p-value среди основных тестов равно $p=0.880$. Проверенные наблюдаемые величины DR1 не показывают повторяющейся направленной оси. Мы считаем это нулевым результатом для DR1 и прекращаем дальнейшее сканирование тех же данных без новых данных, официальных mock-каталогов для свойств или заранее зафиксированного нового семейства наблюдаемых величин.
\end{abstract}

\section{Введение}

Стандартная космологическая модель предполагает статистическую однородность и изотропию на достаточно больших масштабах. В таком случае угловая структура в каталогах галактик и квазаров может возникать из-за космологического кластеринга, маски обзора, эффектов отбора, локальных наблюдательных систематик и статистического шума. Поэтому поиск выделенного направления должен проводиться относительно скорректированной нулевой модели, а не по сырым счетам объектов на небе.

Эта работа задает конкретный наблюдательный вопрос: есть ли в публичных каталогах DESI DR1 признаки слабого дипольного направленного градиента в плотности трассеров или их свойствах после учета основных эффектов обзора и отбора?

Мотивация простая: если выделенное направление существует, оно не обязано проявляться как видимый центр на небе. Оно может входить как слабая статистическая анизотропия в наблюдаемые величины далеких трассеров. Статистическая процедура не предполагает физическую модель такого эффекта. Проверяются только следующие эмпирические гипотезы:

\begin{description}
    \item[$H_0$] После поправок на геометрию обзора, эффекты отбора, random-каталоги и доступные шаблоны систематик проверяемые распределения трассеров DESI DR1 и остаточные наблюдаемые величины совместимы со статистической изотропией.
    \item[$H_1$] Одно или несколько независимых семейств трассеров или свойств показывают устойчивый направленный диполь, который не объясняется геометрией обзора, систематиками или статистическим шумом и повторяется при разбиении по красному смещению, региону или трассеру.
\end{description}

В приведенных ниже тестах такой направленный градиент не обнаружен.

\section{Данные}

\subsection{Каталоги DESI DR1 Large-Scale Structure}

Мы используем публичные clustering-каталоги DESI DR1 Large-Scale Structure (LSS). DESI DR1 включает данные первых 13 месяцев основного обзора, а также единообразно переобработанные данные Survey Validation. В этой работе используются продукты DR1 Iron \texttt{LSScats/v1.5} и соответствующие random-каталоги.

Анализ плотности использует clustering-каталоги QSO в North Galactic Cap (NGC) и South Galactic Cap (SGC), вместе с соответствующими random-каталогами. Random-каталоги здесь принципиальны: сами по себе сырые счета в основном измеряли бы маску обзора и геометрию таргетинга, а не космологическую анизотропию.

\subsection{Value-added каталог FastSpecFit}

Для популяционных и эмиссионных наблюдаемых величин мы используем value-added каталог DESI DR1 FastSpecFit Iron. FastSpecFit моделирует спектры DESI и широкополосную фотометрию, выдавая параметры звездного континуума, эмиссионных линий и производные спектральные величины. В этой работе используются:

\begin{itemize}
    \item \texttt{DN4000\_MODEL} для остатков популяций LRG;
    \item эквивалентные ширины эмиссионных линий QSO, прежде всего \texttt{CIV\_1549\_EW}, \texttt{CIII\_1908\_EW} и \texttt{MGII\_2796\_EW};
    \item столбцы обратной дисперсии для IVAR-взвешенных карт остатков;
    \item диагностические параметры качества подгонки и фотометрические proxy-столбцы как служебные контроли.
\end{itemize}

Наблюдаемая величина QSO \texttt{DN4000\_MODEL} была проверена только как быстрая диагностическая проверка и затем исключена из основного семейства: строгие отборы по качеству оставляют всего 90 строк QSO высокого красного смещения, что делает ее непригодной для научного теста популяционного градиента у таких QSO.

\subsection{Внешние шаблоны систематик}

Для тестов плотности QSO и свойств линий QSO мы включаем внешние HEALPix-шаблоны при \texttt{nside=16}:

\begin{itemize}
    \item прокси галактического поглощения \texttt{EBV};
    \item звездная плотность;
    \item глубина imaging-данных в $g/r/z$;
    \item размер PSF в $g/r/z$;
    \item прокси яркости неба в $g/r/z$.
\end{itemize}

Эти шаблоны используются как регрессионные контроли и не интерпретируются как физические переменные градиента.

\section{Методы}

\subsection{Карты неба и оценка диполя}

Объекты пикселизуются в HEALPix-карты. Если не указано иное, используется \texttt{nside=16}. Для тестов плотности скорректированная карта сверхплотности строится по счетам данных и random-каталогов. Для тестов свойств сначала строятся объектные остатки, а затем они усредняются по небесным пикселям.

Для значения карты $y_p$ в пикселе $p$ и единичного вектора $\hat{n}_p$, направленного в центр пикселя, подгоняется модель
\begin{equation}
    y_p = a_0 + \mathbf{d}\cdot \hat{n}_p + \epsilon_p ,
\end{equation}
где $a_0$ — монопольный член, $\mathbf{d}$ — вектор диполя, а $\epsilon_p$ — остаток. Амплитуда подогнанного диполя равна $|\mathbf{d}|$, а ось диполя переводится в прямое восхождение и склонение.

Подогнанная ось используется как диагностическая величина, а не как самостоятельный результат. Ее интерпретация зависит от распределения при нулевой гипотезе и проверок по трассерам, наблюдаемым величинам, регионам неба и интервалам красного смещения.

\subsection{Остатки популяционных величин и линий}

Для наблюдаемых величин FastSpecFit объектные остатки оцениваются методом взвешенных наименьших квадратов. Общая модель остатков имеет вид
\begin{equation}
    O_i = \beta_0 + \beta_1 z_i + \beta_2 z_i^2 + \sum_j \gamma_j X_{ij} + r_i ,
\end{equation}
где $O_i$ — наблюдаемая величина, $X_{ij}$ — nuisance-контроли, а $r_i$ — остаток, используемый в картах неба.

Для LRG \texttt{DN4000\_MODEL} контролями являются $z$, $z^2$ и \texttt{LOGMSTAR}. Для эквивалентных ширин линий QSO и отношений линий контроли включают $z$, $z^2$, \texttt{LOGMSTAR} при наличии, диагностики подгонки линий \texttt{RCHI2\_LINE} и \texttt{DELTA\_LINECHI2}, фотометрические прокси и прокси светимости, а также перечисленные выше внешние шаблоны систематик.

Карта остатков использует IVAR-взвешенные средние по пикселям:
\begin{equation}
    \bar{r}_p =
    \frac{\sum_{i\in p} r_i w_i}{\sum_{i\in p} w_i},
\end{equation}
где $w_i$ — соответствующая обратная дисперсия.

Для отношений линий вида $\ln(x/y)$ распространяемая обратная дисперсия равна
\begin{equation}
    w_{\ln(x/y)}
    =
    \left[
    \frac{1}{w_x x^2}
    +
    \frac{1}{w_y y^2}
    \right]^{-1}.
\end{equation}
Сохраняются только строки с конечными наблюдаемыми величинами и положительной распространяемой обратной дисперсией.

\subsection{Отборы по качеству}

Основной запуск для LRG \texttt{DN4000\_MODEL} использует:
\begin{itemize}
    \item \texttt{ZWARN == 0};
    \item \texttt{RCHI2\_CONT < 2};
    \item \texttt{DELTACHI2 > 25};
    \item положительный \texttt{DN4000\_IVAR};
    \item винзоризация \texttt{DN4000\_MODEL} по 1-му и 99-му процентилям.
\end{itemize}

Для наблюдаемых величин линий QSO мы не делаем жесткий отбор по \texttt{DELTA\_LINECHI2 > 25}. Проверка выборки QSO высокого красного смещения показала, что такой порог оставил бы только 55 пригодных строк \texttt{CIV\_1549\_EW} и 68 пригодных строк \texttt{CIII\_1908\_EW} из примерно 799,000 сопоставленных QSO высокого красного смещения. Вместо этого \texttt{RCHI2\_LINE} и \texttt{DELTA\_LINECHI2} включаются как служебные контроли, а жесткие отборы оставляют конечные наблюдаемые величины, положительный IVAR и \texttt{ZWARN == 0}.

\subsection{Тесты нулевой гипотезы}

Основной тест нулевой гипотезы для карт свойств — пространственная блочная перестановка:
\begin{enumerate}
    \item разделить небо на грубые HEALPix-блоки с \texttt{block\_nside=2};
    \item переставить паттерны остатков по блокам;
    \item заново подогнать диполь для каждой переставленной карты;
    \item вычислить p-value как долю амплитуд из нулевого распределения, не меньших наблюдаемой амплитуды, используя конечновыборочную поправку $+1$.
\end{enumerate}

Для тестов плотности мы используем пиксельные перестановки и блочные диагностики нулевой гипотезы; интерпретация основана на скорректированных картах после регрессии по random-каталогам и внешним шаблонам.

\subsection{Учет перебора}

Мы сообщаем минимальное p-value по контролируемому наблюдательному семейству и применяем простые поправки Бонферрони и Шидака. Тесты коррелированы, поэтому эти поправки следует читать как консервативные проверки, а не как точное число независимых испытаний. Поскольку все p-values велики, вывод не зависит от конкретной формулы поправки.

\section{Результаты}

\subsection{Остатки популяций LRG}

Основной запуск для LRG проверяет \texttt{DN4000\_MODEL} в регионе NGC при $0.4 \le z < 0.6$.

\begin{table}[h]
\centering
\caption{Остаточный диполь LRG \texttt{DN4000\_MODEL}.}
\begin{tabular}{llrrrrr}
\toprule
Трассер & Регион & Строки & Ампл. & RA & DEC & Блочный p \\
\midrule
LRG & NGC & 342791 & 0.075220 & 247.056 & 38.150 & 0.409182 \\
\bottomrule
\end{tabular}
\end{table}

Этот результат не выделяется на фоне изотропной нулевой модели. Он дает минимальное p-value в контролируемом наблюдательном семействе, но само значение остается далеко внутри нулевого диапазона.

\subsection{Плотность QSO высокого красного смещения}

Тест плотности QSO использует два интервала красного смещения, объединяя NGC+SGC, с DESI random-каталогами и регрессией по внешним шаблонам.

\begin{table}[h]
\centering
\caption{Объединенные тесты плотности QSO после поправки по внешним шаблонам.}
\begin{tabular}{lrrrrr}
\toprule
Диапазон $z$ & $N_{\rm data}$ & Ампл. & RA & DEC & Блочный p \\
\midrule
1.5--2.1 & 432458 & 0.003703 & 87.961 & -78.002 & 0.840432 \\
2.1--3.5 & 366560 & 0.003987 & 329.996 & -4.010 & 0.780044 \\
\bottomrule
\end{tabular}
\end{table}

Тесты с разбиением по регионам также дают высокие p-values.

\begin{table}[h]
\centering
\caption{Тесты плотности QSO с разбиением по регионам.}
\begin{tabular}{llrrrrr}
\toprule
Регион & Диапазон $z$ & $N_{\rm data}$ & Ампл. & RA & DEC & Блочный p \\
\midrule
NGC & 1.5--2.1 & 281438 & 0.009911 & 120.968 & 80.962 & 0.653069 \\
NGC & 2.1--3.5 & 237306 & 0.006319 & 339.286 & 44.405 & 0.876625 \\
SGC & 1.5--2.1 & 151020 & 0.018968 & 187.966 & -54.340 & 0.765847 \\
SGC & 2.1--3.5 & 129254 & 0.017125 & 263.096 & -0.890 & 0.787043 \\
\bottomrule
\end{tabular}
\end{table}

Разбиение по galactic caps не указывает на общую ось, и ни один интервал плотности не дает свидетельства в пользу выделенного направления.

\subsection{Остатки эквивалентных ширин эмиссионных линий QSO}

Объединенные тесты остатков линий QSO высокого красного смещения используют \texttt{CIV\_1549\_EW} и \texttt{CIII\_1908\_EW}.

\begin{table}[h]
\centering
\caption{Объединенные тесты остатков эквивалентных ширин эмиссионных линий QSO.}
\begin{tabular}{lrrrrr}
\toprule
Наблюдаемая величина & Строки & Ампл. & RA & DEC & Блочный p \\
\midrule
\texttt{CIV\_1549\_EW} & 761395 & 0.303061 & 85.281 & 7.464 & 0.958084 \\
\texttt{CIII\_1908\_EW} & 758337 & 0.148250 & 132.129 & -42.987 & 0.904192 \\
\bottomrule
\end{tabular}
\end{table}

Две оси разделены углом 65.953 градуса. Разбиения по региону и красному смещению также не показывают детекции.

\begin{table}[h]
\centering
\caption{Проверки устойчивости остатков линий QSO по региону и красному смещению.}
\begin{tabular}{lllrr}
\toprule
Регион & Диапазон $z$ & Наблюдаемая величина & Строки & Блочный p \\
\midrule
NGC & 1.5--2.1 & \texttt{CIV\_1549\_EW} & 266722 & 0.816367 \\
NGC & 1.5--2.1 & \texttt{CIII\_1908\_EW} & 264221 & 0.972056 \\
NGC & 2.1--3.5 & \texttt{CIV\_1549\_EW} & 228136 & 0.954092 \\
NGC & 2.1--3.5 & \texttt{CIII\_1908\_EW} & 228933 & 0.942116 \\
SGC & 1.5--2.1 & \texttt{CIV\_1549\_EW} & 142895 & 0.988024 \\
SGC & 1.5--2.1 & \texttt{CIII\_1908\_EW} & 141485 & 0.996008 \\
SGC & 2.1--3.5 & \texttt{CIV\_1549\_EW} & 123642 & 0.990020 \\
SGC & 2.1--3.5 & \texttt{CIII\_1908\_EW} & 123698 & 0.966068 \\
\bottomrule
\end{tabular}
\end{table}

\subsection{Предварительно зафиксированное семейство отношений линий QSO}

После того как тесты сырых эквивалентных ширин линий дали нулевые результаты, мы заранее зафиксировали семейство отношений линий перед запуском анализа. Основные наблюдаемые величины:
\begin{align}
    R_{\rm CIV/CIII} &= \ln\left(\frac{\texttt{CIV\_1549\_EW}}{\texttt{CIII\_1908\_EW}}\right), \\
    R_{\rm CIV/MGII} &= \ln\left(\frac{\texttt{CIV\_1549\_EW}}{\texttt{MGII\_2796\_EW}}\right), \\
    R_{\rm CIII/MGII} &= \ln\left(\frac{\texttt{CIII\_1908\_EW}}{\texttt{MGII\_2796\_EW}}\right).
\end{align}

Основное семейство содержит шесть тестов: три отношения линий в двух интервалах красного смещения, NGC+SGC вместе.

\begin{table}[h]
\centering
\caption{Основное предварительно зафиксированное семейство отношений линий QSO.}
\begin{tabular}{llrrrrr}
\toprule
Диапазон $z$ & Наблюдаемая величина & Строки & Ампл. & RA & DEC & Блочный p \\
\midrule
1.5--2.1 & \texttt{LOG\_CIV\_CIII\_EW} & 396908 & 0.005927 & 90.086 & 13.236 & 0.970060 \\
1.5--2.1 & \texttt{LOG\_CIV\_MGII2796\_EW} & 407103 & 0.002747 & 149.263 & -31.179 & 0.994012 \\
1.5--2.1 & \texttt{LOG\_CIII\_MGII2796\_EW} & 403644 & 0.005951 & 251.516 & 25.093 & 0.980040 \\
2.1--3.5 & \texttt{LOG\_CIV\_CIII\_EW} & 345404 & 0.004468 & 171.736 & 10.948 & 0.980040 \\
2.1--3.5 & \texttt{LOG\_CIV\_MGII2796\_EW} & 186086 & 0.017718 & 174.848 & -54.049 & 0.880240 \\
2.1--3.5 & \texttt{LOG\_CIII\_MGII2796\_EW} & 186114 & 0.012246 & 109.881 & 27.207 & 0.980040 \\
\bottomrule
\end{tabular}
\end{table}

Минимальное p-value среди основных тестов равно $p=0.880240$. Для шести основных тестов поправка Бонферрони дает 1.0, а глобальное значение Шидака равно 0.999997. Набор проверок устойчивости по NGC и SGC отдельно имеет минимальное p-value $p=0.724551$. Отношения линий не дают детекции.

\subsection{Сводка с учетом перебора}

Контролируемое наблюдательное семейство, исключая диагностическую проверку QSO \texttt{DN4000\_MODEL}, содержит 17 тестов.

\begin{table}[h]
\centering
\caption{Сводка по основному семейству тестов.}
\begin{tabular}{lrrl}
\toprule
Семейство & Тесты & Минимальный блочный p & Статус \\
\midrule
Остатки популяций LRG & 1 & 0.409182 & нет детекции \\
Скорректированная плотность QSO, вместе & 2 & 0.780044 & нет детекции \\
Плотность QSO по регионам & 4 & 0.653069 & нет детекции \\
Остатки линий QSO, вместе & 2 & 0.904192 & нет детекции \\
Остатки линий QSO по региону/bin & 8 & 0.816367 & нет детекции \\
\bottomrule
\end{tabular}
\end{table}

Минимальное p-value равно $0.4091816367$. Поправка Бонферрони дает 1.0, а глобальное значение Шидака равно 0.999870.

\section{Обсуждение}

Тесты покрывают несколько мест, где мог бы проявиться направленный градиент: плотность QSO высокого красного смещения, остатки звездных популяций LRG, сырые остатки эмиссионных линий QSO, остатки отношений линий QSO, а также отдельные регионы и интервалы красного смещения. Ни один тест не дает малого p-value. Подогнанные оси также не повторяются между независимыми семействами наблюдаемых величин.

Нулевой результат не вызван одной неудачной наблюдаемой величиной. Один и тот же исход получается во всех проверенных семействах DESI DR1.

Ключевой методологический момент состоит в том, что анализ не опирается на заявления об анизотропии сырых счетов. Тесты плотности корректируются с помощью random-каталогов и внешних шаблонов. Тесты свойств перед проекцией на небо очищаются по красному смещению, диагностикам качества подгонки, прокси светимости и imaging-шаблонам. P-values получены из пространственных блочных перестановок, которые лучше сохраняют крупные корреляции на небе, чем независимые перестановки пикселей.

Это не доказывает точную космическую изотропию. Результат говорит только о том, что проверенные здесь сигнатуры направленного градиента не обнаружены в текущих публичных продуктах DESI DR1 при использованных поправках.

\section{Ограничения}

Для тестов остатков свойств пока нет официальных mock-каталогов, откалиброванных под точные остатки наблюдаемых величин FastSpecFit. Тесты плотности лучше закреплены DESI random-каталогами, но остатки свойств зависят от адекватности модели остатков и шаблонов систематик.

Основной тест популяций LRG сейчас покрывает один интервал красного смещения и регион NGC. Дополнительные тесты популяций LRG и ELG могли бы быть полезны, но дальнейшее сканирование DR1 без предварительной фиксации повышало бы риск подбора результата после просмотра данных.

QSO \texttt{DN4000\_MODEL} не полезен на высоком красном смещении при использованных строгих отборах по качеству. Он был исключен из основного семейства после диагностического запуска, в котором осталось только 90 строк.

Все поправки на множественное тестирование являются приближенными, поскольку связанные разбиения неба и наблюдаемые величины коррелированы. Так как все p-values велики, вывод не чувствителен к точной форме поправки.

Наконец, это первичный поиск направленного градиента, а не полный космологический параметрический анализ. Здесь не подгоняется физическая анизотропная космология, и эта работа не заменяет стандартные анализы DESI BAO или искажений в пространстве красных смещений.

\section{Условие остановки}

После того как предварительно зафиксированное семейство отношений линий QSO также не дало детекции, мы установили условие остановки для сканирования DESI DR1. Дальнейшие поиски на тех же данных DR1 прекращаются, если не выполняется одно из следующих условий:
\begin{itemize}
    \item используется новый релиз данных DESI;
    \item становятся доступны официальные mock-каталоги для точного оценивателя свойств и остатков;
    \item документированная методологическая ошибка делает один из существующих результатов недействительным;
    \item новый набор внешних шаблонов систематик существенно меняет модель поправок;
    \item новое физическое семейство наблюдаемых величин зафиксировано до запуска анализа.
\end{itemize}

Это условие остановки нужно, чтобы предотвратить неконтролируемые поиски наблюдаемых величин после просмотра результатов и повторных нулевых исходов.

\section{Выводы}

Мы искали направленные градиенты в DESI DR1 LSS и наблюдаемых величинах FastSpecFit с помощью карт плотности, остатков популяций галактик, остатков эмиссионных линий QSO и предварительно зафиксированных отношений линий QSO. Тесты используют random-каталоги, внешние шаблоны систематик, моделирование остатков, HEALPix-подгонку диполя и пространственные блочные p-values для нулевой гипотезы.

Направленный градиент не обнаружен.

Эмпирический статус этого поиска по DR1:
\begin{quote}
DESI DR1 LSS + FastSpecFit: поиск направленного градиента, нет детекции.
\end{quote}

Дальнейшая работа должна быть сосредоточена на проверке по DESI DR2/DR3, официальных mock-каталогах для остатков свойств и заранее зафиксированных семействах наблюдаемых величин, а не на дополнительных незарегистрированных сканированиях тех же данных DR1.

\section*{Доступность данных и кода}

Анализ был выполнен локально в Python-проекте \texttt{cosmo\_genesis\_gradient}. Ключевые сгенерированные артефакты:
\begin{itemize}
    \item \texttt{outputs/reports/observational\_gradient\_master\_null\_report.md};
    \item \texttt{outputs/tables/observational\_gradient\_master\_null\_tests.csv};
    \item \texttt{outputs/reports/qso\_line\_ratio\_preregistration.md};
    \item \texttt{outputs/reports/qso\_line\_ratio\_preregistered\_results.md};
    \item \texttt{outputs/tables/qso\_line\_ratio\_preregistered\_results.csv};
    \item \texttt{outputs/reports/observational\_stop\_condition.md}.
\end{itemize}

Код организован как Python-проект с command-line entry points под \texttt{cosmo-gradient}.

\begin{thebibliography}{9}

\bibitem{desidr1docs}
DESI Collaboration,
\newblock документация DESI Data Release 1,
\newblock \url{https://data.desi.lbl.gov/doc/releases/dr1/}.

\bibitem{desidr1paper}
DESI Collaboration,
\newblock Data Release 1 of the Dark Energy Spectroscopic Instrument,
\newblock arXiv:2503.14745,
\newblock \url{https://arxiv.org/abs/2503.14745}.

\bibitem{desilss}
DESI Collaboration,
\newblock каталог DESI DR1 LSS,
\newblock \url{https://data.desi.lbl.gov/public/dr1/survey/catalogs/dr1/LSS/iron/LSScats/v1.5/}.

\bibitem{desiorg}
DESI Collaboration,
\newblock документация по организации данных DESI,
\newblock \url{https://data.desi.lbl.gov/doc/organization/}.

\bibitem{fastspecfit}
J. Moustakas et al.,
\newblock документация FastSpecFit,
\newblock \url{https://fastspecfit.readthedocs.io/en/2.1.2/}.

\bibitem{desiagnqso}
DESI Collaboration,
\newblock документация DESI DR1 AGN/QSO VAC,
\newblock \url{https://data.desi.lbl.gov/doc/releases/dr1/vac/agnqso/}.

\end{thebibliography}

\end{document}
